\documentclass[letterpaper, paper,11pt]{AAS}

\usepackage{bm}
\usepackage{amsmath}
\usepackage{subfigure}
\usepackage[colorlinks=true, pdfstartview=FitV, linkcolor=black, citecolor= black, urlcolor= black]{hyperref}
\usepackage{overcite}
\usepackage{footnpag}
\usepackage{tikz}
\usepackage{float}
\usepackage{comment}
\usepackage{booktabs}
\usetikzlibrary{shapes.geometric, arrows.meta, positioning}
\tikzstyle{process} = [rectangle, minimum width=3.5cm, minimum height=1cm, text centered, draw=black, fill=blue!10]
\tikzstyle{decision} = [diamond, minimum width=3cm, minimum height=1cm, text centered, draw=black, fill=yellow!20]
\tikzstyle{startstop} = [ellipse, minimum width=3cm, text centered, draw=black, fill=red!20]
\tikzstyle{arrow} = [thick, ->, >=stealth]
\PaperNumber{25-787}

\begin{document}

% ==============================================================================
% SECTION: METHODOLOGY
% ==============================================================================
\section{Methodology}
\label{sec:methodology}

The architecture proposed in this study introduces a hybrid Symplectic-AI framework designed to overcome the computational bottlenecks inherent in Next-Generation Space Traffic Management. Traditional numerical propagators face exponential scaling challenges when applied to fragmentation events involving thousands of objects. To address this, the methodology is divided into three sequential phases: physics-informed surrogate modeling, sequence-based spatial filtering, and autonomous event reconstruction via a high-speed memory bridge.

% ------------------------------------------------------------------------------
% Phase 1: PINNs
% ------------------------------------------------------------------------------
\subsection{Phase 1: Physics-Informed Surrogate Modeling}
\label{subsec:method_phase1}

The standard calculation of atmospheric drag for thousands of objects in Low Earth Orbit (LEO) requires continuous evaluations of localized atmospheric density. To bypass the computational latency of these standard exponential decay models, a Physics-Informed Neural Network (PINN) is introduced as a high-speed surrogate model.

The baseline thermospheric density, $\rho(h)$, at altitude $h$, is analytically approximated by:
$$ \rho(h) = \rho_0 \exp\left(-\frac{h - h_0}{H}\right) $$
where $\rho_0$ is the reference density at reference altitude $h_0$, and $H$ is the atmospheric scale height. Standard deep neural networks are prone to non-physical interpolation when approximating this function. Therefore, the PINN is constrained by the analytical derivative of the physical system:
$$ \frac{d\rho}{dh} = -\frac{\rho(h)}{H} $$

The neural network is trained using a composite loss function that penalizes both data-driven mean squared error ($\mathcal{L}_{Data}$) and violations of the underlying physical gradient ($\mathcal{L}_{Physics}$):
$$ \mathcal{L}_{Total} = \frac{1}{N_{data}} \sum_{i=1}^{N_{data}} (\hat{\rho}_i - \rho_i)^2 + \lambda \frac{1}{N_{colloc}} \sum_{j=1}^{N_{colloc}} \left( \frac{\partial \hat{\rho}_j}{\partial h} + \frac{\hat{\rho}_j}{H \ln(10)} \right)^2 $$

By enforcing this gradient penalty ($\lambda$), the surrogate model ensures strictly monotonic, physically valid density profiles across the LEO regime, enabling rapid tensor-based evaluations for downstream perturbations.

% ------------------------------------------------------------------------------
% Phase 2: KD-Tree & LSTM
% ------------------------------------------------------------------------------
\subsection{Phase 2: Spatial Filtering and Sequence-Based Conjunction Prediction}
\label{subsec:method_phase2}

When a fragmentation event generates a localized cloud of thousands of objects, evaluating pairwise conjunctions presents an $\mathcal{O}(N^2)$ computational complexity. To achieve scalability, the architecture implements a spatial filtering protocol followed by a sequence classification model.

First, the Cartesian state vectors generated by the Symplectic library are ingested into a $K$-Dimensional Tree (KD-Tree). The KD-Tree partitions the operational volume, isolating regional neighborhoods and reducing the conjunction search space complexity to $\mathcal{O}(N \log N)$. 

Once the boundary-case objects are isolated, their relative state histories are converted into sliding time-series sequences. A Long Short-Term Memory (LSTM) recurrent neural network is deployed to classify these temporal sequences. The network takes an input matrix $\mathbf{X}_t \in \mathbb{R}^{S \times 6}$, where $S$ is the sequence length, representing the relative position and velocity vectors:
$$ \mathbf{x}_i = [\Delta x, \Delta y, \Delta z, \Delta v_x, \Delta v_y, \Delta v_z]^T $$

The LSTM architecture identifies the non-linear kinematic signatures of decaying orbits and predicts a binary output indicating whether the relative distance will breach a predefined safety threshold within the prediction horizon.

% ------------------------------------------------------------------------------
% Phase 3: RL & Memory Bridge
% ------------------------------------------------------------------------------
\subsection{Phase 3: Autonomous Event Reconstruction via Reinforcement Learning}
\label{subsec:method_phase3}

The final phase addresses the inverse problem of event reconstruction: determining the origin state (time, position, and velocity distribution) of a debris-generating event based on sparse observational data. This is formulated as a Markov Decision Process (MDP) and solved using Proximal Policy Optimization (PPO).

The RL agent interacts with an operational environment parameterized by a continuous action space $\mathbf{a}_t \in \mathbb{R}^5$:
$$ \mathbf{a}_t = [\delta t, \delta x, \delta y, \delta z, \delta v_{spread}] $$

To evaluate the fitness of each guessed action, the environment must propagate a simulated debris cloud forward in time to match the observation epoch. Traditional file-based Input/Output (I/O) methods create a prohibitive latency barrier when executing the millions of integration loops required for RL convergence. 

To resolve this, the C-based Symplectic library is compiled as a shared dynamic library (\texttt{.so}) and linked directly to the Python-based Gymnasium environment via foreign function interfaces (\texttt{ctypes}). This memory bridge allows the RL agent to allocate simulated particles, execute fixed-timestep symplectic integrations, and retrieve the final state vectors entirely within hardware RAM.

The state observation provided to the agent is the error between the statistical moments of the simulated cloud and the true observational data. The agent is driven by a reward function targeting the Root Mean Square Error (RMSE) of the spatial distribution:
$$ R_t = - \frac{1}{\gamma} \sqrt{ ||\boldsymbol{\mu}_{sim} - \boldsymbol{\mu}_{obs}||^2 + ||\boldsymbol{\sigma}_{sim} - \boldsymbol{\sigma}_{obs}||^2 } $$
where $\gamma$ is a standardizing scaling factor, $\boldsymbol{\mu}$ represents the volumetric centroid of the debris cloud, and $\boldsymbol{\sigma}$ represents the spatial dispersion. The agent is strictly bounded to physically valid Cartesian limits to prevent gravitational singularities during integration, ensuring stable gradient updates across parallelized rollout workers.

\end{document}